\documentclass[a4paper, serif, 10pt]{moderncv}

%% moderncv
% color and style
\moderncvcolor{black}
\moderncvstyle{banking}

% renew moderncv command to adapt for CJK colon
\renewcommand*{\cvitem}[3][.25em]{%
  \ifstrempty{#2}{}{\hintstyle{#2}：}{#3}%
  \par\addvspace{#1}}

\renewcommand*{\cvitemwithcomment}[4][.25em]{%
  \savebox{\cvitemwithcommentmainbox}{\ifstrempty{#2}{}{\hintstyle{#2}：}#3}%
  \setlength{\cvitemwithcommentmainlength}{\widthof{\usebox{\cvitemwithcommentmainbox}}}%
  \setlength{\cvitemwithcommentcommentlength}{\maincolumnwidth-\separatorcolumnwidth-\cvitemwithcommentmainlength}%
  \begin{minipage}[t]{\cvitemwithcommentmainlength}\usebox{\cvitemwithcommentmainbox}\end{minipage}%
  \hfill% fill of \separatorcolumnwidth
  \begin{minipage}[t]{\cvitemwithcommentcommentlength}\raggedleft\small\itshape#4\end{minipage}%
  \par\addvspace{#1}}

%% page layout/margins
\usepackage[top=2.5cm, bottom=2.5cm, left=1.5cm, right=1.5cm]{geometry}

\nopagenumbers{}

%% Babel config for English language
\usepackage[english]{babel}

%% fontspec
\usepackage{fontspec}

\IfFontExistsTF{Linux Libertine}{
  \setmainfont[Ligatures={TeX, Common}, Numbers=Lining, ItalicFont=Linux Libertine]{Linux Libertine}
}{}
\IfFontExistsTF{Linux Libertine O}{
  \setmainfont[Ligatures={TeX, Common}, Numbers=Lining, ItalicFont=Linux Libertine O]{Linux Libertine O}
}{}

%% CTeX
% CJK support, used to show CJK characters in the resume
%
% - heading=false: disable ctex heading
% - punct=kaiming: use kaiming punctuations styles for CJK
% - scheme=plain: use plain scheme, do not override \normalsize font size
% - space=auto: space settings for CJK characters
%
% CJK fonts are resolved per locale language, e.g. Noto Serif/Sans CJK JP
% for Japanese resumes, with fallback to the next region if not installed
%
% ref:
% - http://ctan.mirrorcatalogs.com/language/chinese/ctex/ctex.pdf
\usepackage[UTF8, heading=false, punct=kaiming, scheme=plain, space=auto]{ctex}

\IfFontExistsTF{Noto Serif CJK TC}{
  \setCJKmainfont{Noto Serif CJK TC}
}{}
\IfFontExistsTF{Noto Sans CJK TC}{
  \setCJKsansfont{Noto Sans CJK TC}
}{}

%% line spacing
\usepackage{setspace}
\setstretch{1.125}

%% URL styling - use normal text instead of monospace
\urlstyle{same}

%% Patch \httplink and \httpslink to handle full URLs
% The original moderncv macros always prepend http:// or https:// to URLs.
% This causes issues when the URL already has a protocol (e.g., https://example.com)
% resulting in malformed URLs like https://https://example.com.
% This patch checks if the URL already contains "://" and uses it directly if so.
% Note: We use \str_set:Nx (x-expansion) to expand macros like \@homepage first.
\ExplSyntaxOn
\str_new:N \g_yamlresume_url_str

\RenewDocumentCommand{\httpslink}{O{}m}{
  \str_set:Nx \g_yamlresume_url_str {#2}
  \str_if_in:NnTF \g_yamlresume_url_str {://}
    { \url{#2} }
    { \url{https://#2} }
}

\RenewDocumentCommand{\httplink}{O{}m}{
  \str_set:Nx \g_yamlresume_url_str {#2}
  \str_if_in:NnTF \g_yamlresume_url_str {://}
    { \url{#2} }
    { \url{http://#2} }
}
\ExplSyntaxOff

\name{林怡君}{}
\title{資深商業分析師 | 數據驅動決策專家}
\phone[mobile]{+886 912 345 678}
\email{yichun.lin@example.com}
\homepage{https://www.linkedin.com/in/yichunlin}

\address{中國臺灣，信義區，台北市，台北市信義區松仁路100號，110}{}{}

\extrainfo{{\small \faLinkedin}\ \href{https://www.linkedin.com/in/yichunlin}{@yichunlin} {} {} {} • {} {} {} 
{\small \faGithub}\ \href{https://github.com/yichunlin}{@yichunlin}}

\begin{document}

\maketitle

\section{簡介}

\cvline{}{\begin{itemize}
\item 具備 8 年以上商業分析經驗，專精於金融與電信產業的數據驅動決策
\item 擅長需求訪談、流程優化與商業案例撰寫，成功協助企業提升營運效率達 20\%
\item 熟練運用 SQL、Python、Tableau 等工具進行數據分析與視覺化
\item 擁有卓越的跨部門溝通與專案管理能力，能有效協調利害關係人達成目標
\end{itemize}}
\section{教育背景}

\cventry{2010年9月–2014年6月}
        {學士，企業管理學系，成績：3.9}
        {國立臺灣大學}
        {\url{https://www.ntu.edu.tw}}
        {}
        {\begin{itemize}
\item 主修企業管理，輔修資訊管理，奠定商業分析與數據處理的紮實基礎
\item 畢業專題「電子商務消費者行為分析」獲得系上最佳專題獎
\item 積極參與商業競賽，獲得全國大專院校商業個案大賽優勝
\end{itemize}
\textbf{課程}：統計學、作業研究、財務管理、行銷管理、資訊管理、商業決策分析}
\section{工作經歷}

\cventry{2020年1月至今}
        {資深商業分析師}
        {富邦金融控股股份有限公司}
        {\url{https://www.fubon.com}}
        {}
        {\begin{itemize}
\item 領導 5 人分析團隊，負責銀行與保險業務的數據分析與流程優化
\item 主導智能客服導入專案，整合 NLP 技術，降低客服成本 30\% 並提升客戶滿意度 15\%
\item 建立客戶流失預測模型，準確率達 85\%，協助挽留高價值客戶，增加年營收 5,000 萬元
\item 與 IT、業務、風控等部門協作，定義需求規格並確保專案如期交付
\item 定期向高階管理層彙報分析結果與策略建議，推動數據驅動文化
\end{itemize}
\textbf{關鍵字}：需求分析、流程優化、數據建模、利害關係人管理、金融科技}

\cventry{2016年7月–2019年12月}
        {商業分析師}
        {台灣大哥大股份有限公司}
        {\url{https://www.taiwanmobile.com}}
        {}
        {\begin{itemize}
\item 負責電信資費方案與用戶行為分析，提供定價策略建議，提升 ARPU 8\%
\item 開發 Tableau 儀表板，即時監控關鍵績效指標，縮短報表產出時間 50\%
\item 執行客戶分群與生命周期管理，設計精準行銷活動，提升轉換率 12\%
\item 參與跨部門敏捷開發團隊，擔任產品負責人與開發團隊的橋樑
\end{itemize}
\textbf{關鍵字}：電信分析、Tableau、客戶分群、定價策略、敏捷開發}

\cventry{2014年9月–2016年6月}
        {初級商業分析師}
        {勤業眾信聯合會計師事務所}
        {\url{https://www2.deloitte.com/tw}}
        {}
        {\begin{itemize}
\item 協助金融業客戶進行流程梳理與系統需求訪談，撰寫商業需求文件
\item 參與核心銀行系統升級專案，負責測試案例設計與使用者驗收測試
\item 運用 SQL 進行數據驗證與差異分析，確保資料遷移正確性
\item 獲得年度最佳新人獎，並提前晉升為商業分析師
\end{itemize}
\textbf{關鍵字}：商業需求文件、流程梳理、系統測試、SQL、銀行系統}
\section{語言}

\cvline{中文}{母語或雙語水平 \hfill \textbf{關鍵字}：母語}
\cvline{英語}{完全專業水平 \hfill \textbf{關鍵字}：TOEIC 950、IELTS 7.5}
\section{專業技能}

\cvline{商業分析}{專家 \hfill \textbf{關鍵字}：需求訪談、流程優化、商業案例、利害關係人管理、使用者故事}
\cvline{數據分析}{高級 \hfill \textbf{關鍵字}：SQL、Python、Tableau、Excel、R}
\cvline{專案管理}{高級 \hfill \textbf{關鍵字}：Agile、Scrum、JIRA、風險管理、時程規劃}
\cvline{財務分析}{中級 \hfill \textbf{關鍵字}：財務建模、預算規劃、ROI分析、成本效益分析}
\section{榮譽}

\cventry{2022年12月}
        {富邦金融控股股份有限公司}
        {年度最佳員工}
        {}
        {}
        {表彰在智能客服導入專案中展現卓越領導力與分析能力，成功降低營運成本並提升客戶體驗。}

\cventry{2019年11月}
        {台灣大哥大股份有限公司}
        {傑出專案獎}
        {}
        {}
        {因客戶流失預測模型專案成效顯著，有效提升客戶留存率，獲頒年度傑出專案獎。}
\section{證書}

\cventry{2021年3月}
        {IIBA}
        {Certified Business Analysis Professional (CBAP)}
        {\url{https://www.iiba.org/certification/cbap}}
        {}
        {}

\cventry{2018年6月}
        {PMI}
        {Project Management Professional (PMP)}
        {\url{https://www.pmi.org/certifications/project-management-pmp}}
        {}
        {}
\section{出版刊物}

\cventry{2021年9月}
        {數據驅動的商業決策：台灣金融業的應用}
        {台灣管理學刊}
        {\url{https://example.com/journal/article}}
        {}
        {\begin{itemize}
\item 探討金融業如何運用大數據與機器學習優化商業決策流程
\item 分析台灣多家銀行的實際案例，提出數據治理與分析框架建議
\item 研究結果顯示數據驅動文化可顯著提升企業競爭力
\end{itemize}}
\section{項目}

\cventry{2021年3月–2021年11月}
        {為富邦銀行導入AI智能客服系統，提升客戶服務效率與品質}
        {智能客服導入專案}
        {\url{https://example.com/project/ai-customer-service}}
        {}
        {\begin{itemize}
\item 主導需求訪談與流程設計，整合NLP技術與現有客服系統
\item 協調IT團隊與業務單位，確保系統符合使用者需求
\item 上線後客服成本降低30\%，客戶滿意度提升15\%
\end{itemize}
\textbf{關鍵字}：智能客服、NLP、流程設計、專案管理}

\cventry{2020年1月–2020年8月}
        {建立機器學習模型預測客戶流失風險，協助制定挽留策略}
        {客戶流失預測模型}
        {\url{https://example.com/project/churn-prediction}}
        {}
        {\begin{itemize}
\item 使用Python與Scikit-learn建立分類模型，準確率達85\%
\item 分析客戶行為特徵，找出關鍵流失因子
\item 協助業務團隊設計精準挽留方案，增加年營收5,000萬元
\end{itemize}
\textbf{關鍵字}：機器學習、客戶流失、Python、預測模型}
\section{興趣愛好}

\cvline{閱讀}{商業策略、數據科學、歷史}
\cvline{運動}{馬拉松、游泳、登山}
\section{志願服務}

\cventry{2022年3月至今}
        {志工講師}
        {台灣商業分析師協會}
        {\url{https://www.tbaa.org.tw}}
        {}
        {\begin{itemize}
\item 定期舉辦商業分析工作坊，分享數據驅動決策的實務經驗
\item 指導年輕分析師進行專案規劃與職涯發展
\item 協助協會推廣商業分析專業認證與最佳實踐
\end{itemize}}
    
\end{document}